\begin{abstract}%
We study how the gap between latent and intrinsic dimensionality
shapes memorization in diffusion models. Extending recent spectral
analysis to anisotropic latent data, we find that excess null
dimensions populate a new ``noise-dimension'' bulk in the score
network's feature-correlation spectrum, sitting between
population-driven and sample-specific eigenmodes. Because gradient
flow learns modes in decreasing-eigenvalue order, this bulk acts as
a buffer: each excess latent direction adds a mode that must be
absorbed before sample-specific memorization begins. Synthetic
random-feature and MLP experiments isolate this mechanism. On
CelebA and CIFAR-10 latent diffusion, memorization is
non-monotone in $\dlat$: it first rises as the VAE becomes
expressive enough to reconstruct the training subset, then falls at
larger latent widths while FID is best at intermediate dimensions.
Finally, a data-spectrum surrogate computed from the frozen VAE
encoder predicts much of the empirical memorization trajectory,
suggesting a route from latent geometry to stopping-time estimates.
\end{abstract}
